\chapter{Discussion}
\label{ch:discussion}

\section{Answers to the research questions}

\paragraph{\RQ1: does in-context learning degrade under shift?}
Yes, and concept shift, realised by naming, causes the most degradation. We also saw that reversing the shortcut costs about half as much, and unevenly: where the model's numeric prior agrees with the shortcut the cost is large, and where it opposes it there is almost none, because the prior reads the shortcut in the query with or without demonstrations. Averaged over tasks the prior's part cancels, and what remains is the shortcut the demonstrations taught; hiding it from the demonstrations shrinks the cost in every strategy and naming.

\begin{figure}[H]
\centering
\includegraphics[width=0.55\textwidth]{\figpath/fig_rq2_interaction.pdf}
\caption{The one \RQ2 contrast that survives Holm correction (abstract names). Each grey line is one task; blue is the mean with a 95\% interval (hierarchical bootstrap, 24 tasks). Neutralising the shortcut with \strat{counter\_spurious} costs \AUROC under covariate shift and roughly breaks even under its reversal; the rise between the two, $+0.11$ (Holm $p = 0.024$, 17 of 24 tasks), is the interaction.}
\label{fig:rq2-interaction}
\end{figure}

\newpage
\paragraph{\RQ2: do different strategies help under different shifts?}
Only weakly. The strategies differ far less than the namings do (\autoref{fig:strategies-by-naming}), and none lifts Qwen above chance under flipped names, including the one designed for prior conflict (+0.023 over \strat{label\_diversity}, not significant). Of the five planned \RQ2 contrasts, only one survives Holm correction (\autoref{fig:rq2-interaction}): \strat{counter\_spurious} fares better against \strat{label\_diversity} under shortcut reversal than under covariate shift, mainly because it costs \AUROC under covariate shift; its gain under reversal itself is not significant.

\paragraph{\RQ3: does the model rely on the right features, and why does the prior win?}

Under abstract names it tends to rely on the shortcut, whereas under aligned names it relies on the rule features, and under flipped names on the rule features but in the direction the names suggest. Its self-reports do not match the features it uses. The internal evidence fits the prior winning: the prior is present in the model's state from layer 2, while the demonstrated rule barely appears there, and the answer forms in the last six layers, where the names decide it (\autoref{ch:mechanism}).

\paragraph{\RQ4: does a learned selector help?}
Not on the instruction-tuned model, here Qwen2.5-7B-Instruct: none of the six pre-registered contrasts favours \sata, and taking its top-$k$ rows instead of drawing does no better. On the stand-in learner it is the only method with no weak environment, and in an exploratory comparison it helps the base model (\autoref{app:base}).

\section{The results in context}
\citet{wei2023larger} found that overriding semantic priors with in-context evidence emerges with scale. Our results are consistent with that picture and sharpen it for Qwen2.5-7B-Instruct and its base model. The instruction-tuned model is held more tightly by its priors (the base model loses less to flipped names, and a selector helps it). Although two models cannot show that instruction tuning is the cause, in what we tested the cost of conflict depended more on the names than on which evidence (demonstrations) was chosen. The graded picture of \citet{krishnakumar2025semantic} and \citet{kossen2024incontext}, in which mid-sized models follow flipped labels part of the time, fits the fact that demonstration labels still matter here (gold minus shuffled labels: +0.06 to +0.11 \AUROC, as seen in \autoref{ch:shift}) without being able to overturn the names.

\citet{halawi2024overthinking} find that under false demonstrations, the later layers tend to switch to the demonstrated labels. We find a late critical layer as well, but the switch it marks is driven mainly by the names, as the zero-shot runs show: in our setting the later layers are where the semantic priors are applied, not where the demonstrations take over. \citet{minder2025controllable} report a single context-versus-prior direction in Llama-3.1-8B but, for our numeric task, a mean-difference direction of that kind was a bias rather than a switch.

\citet{nastl2024causal} found that restricting predictors to causal features did not improve generalisation on real tabular shifts. We found the same for demonstrations without the shortcut on ID queries, although we also found a clear gain under reversal. Finally, learned selectors trained on model feedback \citep{rubin2022learning, ye2023compositional} are usually evaluated in distribution and on tasks the model already knows; our results, on the other hand, suggest that a selector's value depends on how much the model reads the evidence at all.

\section{Anticipated criticism}
\label{sec:criticism}

\subsection*{\RQ1}

% \begin{figure}[H]
% \centering
% \includegraphics[width=0.85\textwidth]{\figpath/.pdf}
% \label{fig:rq2-interaction}
% \end{figure}

\critique{Covariate shift showing no effect means your covariate shift is too weak.}

\critique{Naming is not concept shift; $P(y \mid x)$ never changes.}
{From the model's side it does: the same numbers under antonym names imply the opposite mapping, which is what concept shift means for a learner whose prior is semantic. We use the term in that operational sense and say so (\autoref{sec:environments}).}


\subsection*{\RQ2}

\critique{The strategies are weak straw men.}

Two observations suggest they are not. First, each follows a published idea (\autoref{tab:strategies}), and two of them, \strat{rule\_diversity} and \strat{counter\_spurious}, read the ground truth, which no deployed selector could. They are best-case versions of their ideas, so if they cannot beat balanced random, a realistic version is unlikely to, at least at $k = 8$. Second, the strategies are not inert: they change Qwen's behaviour the way their ideas predict. \strat{counter\_spurious}, which makes the shortcut uninformative, all but removes the drop from ID to spurious reversal under every naming: in \autoref{fig:strategies-by-naming} its ID and reversal dots nearly coincide in each panel (gaps of $-0.03$ to $0.01$, against $0.05$ to $0.07$ for \strat{label\_diversity}). It pays for this with lower ID \AUROC, so it does not beat \strat{label\_diversity} overall. The same holds for the others: no strategy reliably beats \strat{label\_diversity}, a random draw with $k/2$ rows of each label. Across six strategies, three namings and three environments, the typical difference from it is 0.025 \AUROC; 40 of the 54 differences have intervals that include zero, 10 lie below it and 4 above, and no strategy is clearly ahead in more than one environment under any naming. Each balanced strategy shows the same label counts as \strat{label\_diversity}, so these comparisons separate which rows are chosen from how many of each label are shown, which biases predictions on its own \citep{zhao2021calibrate}. Choosing rows by a selection idea therefore adds little over picking them at random with the labels balanced, and what it gains in one environment it tends to give back in another. The criticism does apply to selectors we did not test, such as retrievers trained on the model's own feedback; \RQ4 tests one learned selector.

\critique{24 tasks are too few to detect anything.}

They are too few to detect small effects, and we do not claim otherwise. The intervals for the strategy contrasts are about 0.1 \AUROC wide, so effects below about 0.05 would usually be missed; halving that width would take roughly four times as many tasks (assuming that the interval width shrinks with $\sqrt{n_T}$, where $n_T$ is the number of tasks). The design does detect effects of about 0.1: the cost of flipped names (0.20) and the \strat{counter\_spurious} interaction (0.11) both survive Holm correction, and \strat{feature\_range}'s loss under covariate shift ($-0.08$) has an interval entirely below zero (\autoref{tab:contrasts}). What the null results rule out is a large benefit. For the four strategies tested against \strat{label\_diversity} under their target shifts, the largest gain inside any interval is 0.09 \AUROC, under half the cost of flipped names, and two of the four exclude any gain above 0.05; the six \sata contrasts give the same ceiling (\autoref{ch:sata}). A strategy that helps by a few hundredths of \AUROC may exist at $k = 8$ on this model; one that helps as much as the names hurt is ruled out though.

\newpage

\subsection*{\RQ3}

\critique{The logit lens is unreliable for early layers.}

There is a possibility that it can be biased for early layers \citep{belrose2023eliciting}.
The final layer matches the model's output exactly, and the linear probes avoid that bias. The probes tell the same story: the prior is active early in the model, and the rule stays weak throughout.

\critique{Steering failing proves nothing about steering in general.}

Agreed. We tested one family of directions (a label-averaged mean difference at five layers, applied to the query) on one 7B model with eight numeric demonstrations, within a small budget. We claim only that this direction acts as a bias at this scale, because steering adds the same vector $v_l$ to every query, whatever its label, and we explain why it must: following the demonstrations means opposite things for the two labels. When the demonstrations imply label 0, following them means pushing the margin down; when they imply label 1, it means pushing it up, and the measured cosine between the two per-label directions was approximately $-0.9$, i.e.\ they point in nearly opposite directions. Within-cell \AUROC (a ``cell'' is our term for one fixed prompt setup: one task, one strategy, one demonstration seed, one naming and one environment) cannot see a push that moves every margin together, which is why \autoref{fig:steer} reports the margins as well: at layer 24 they fall by about 9 logits, towards answering 0, while \AUROC stays flat, so the direction shifts every answer at once rather than separating the labels.

\critique{Self-reported importance is a known failure; why test it?}

Because it is the cheapest explanation a practitioner would ask for, and because the comparison with behavioural reliance under each naming shows that the self-report follows the names, not the model's behaviour.

\newpage
\subsection*{\RQ4}
\critique{\sata failed because it was badly trained.}

It reached its own oracle's level on the stand-in learner and was the only method with no weak environment there; the base model benefits from the same checkpoint; and ranking instead of drawing does no better. The results point instead to what the instruction-tuned model does with the evidence it is given, as discussed earlier.

\critique{The stand-in learner is a poor proxy for the model, so the selector was optimised for the wrong thing.}

This is the main weakness of the selector chapter, and we measured it. On 8 held-out tasks shown to Qwen, the stand-in learner ranks the strategies across environments much as Qwen does under abstract names: we saw the correlation to be about 0.7, but only weakly under named features, about 0.5. Two results suggest that the proxy still targeted something real. \sata's demonstrations carry the rule, because the stand-in learner learns from them in every environment, and the base model also uses them: with the same checkpoint, its weakest environment is better than any other strategy's (exploratory; \autoref{app:base}). What the proxy did miss, however, is how little the instruction-tuned model uses the evidence once the names and its prior pull the other way. So the direct fix may be to train the selector on Qwen's own answers rather than on a stand-in: running Qwen with many random demonstration sets and scoring each row by how much it helps, just as \citet{chang2023data} do to find helpful demonstrations. However, the inherent bias of a model of this scale would also affect the demonstrations chosen by such a selector, so a larger model, whether in the Qwen family or another, would be more appropriate. Although this was planned, it could not be run within our budget.


\newpage
\section{Limitations}
\label{sec:limitations}

Every negative result in this thesis holds for the conditions we could test, and those conditions were narrow.

\begin{itemize}[leftmargin=1.5em]
  \item \textbf{Scale.} One model family at 7B: Qwen2.5-7B-Instruct and, exploratorily, its base model. Due to time and budget limitations, Llama-3.1-8B-Instruct was dropped when it would not answer in format, and larger models, which \citet{wei2023larger} find follow contradicting evidence better, were out of reach.
  \item \textbf{Demonstrations.} $k = 8$, chosen on the pilot because $k = 32$ did not pass the learnability check. More demonstrations might give the rule more weight.
  \item \textbf{Tasks.} Synthetic numeric tasks with three load-bearing features, two domains and 24 evaluation tasks. Real tabular data, text features or other rule families may behave differently.
\item \textbf{Metric.} Within-cell \AUROC cannot see a shift that moves every query's margin by the same amount; balanced accuracy, its alternative, was distorted by calibration.
  \item \textbf{Proxies.} The selector was trained against a stand-in learner that tracks the model only moderately under named features.
\item \textbf{Budget.} The whole project used a few tens of US dollars of rented GPU time, and the later stages were sized to fit. Steering tested one kind of direction (a single averaged direction, added to the query). Two selectors that learn from Qwen's own answers rather than from a stand-in were planned but not run, because each needs many runs of Qwen for every training task: one that scores each candidate row by how much it helps Qwen \citep{chang2023data}, and one trained by reinforcement learning to add demonstrations one at a time \citep{zhang2022active}. They are the most direct test of whether the stand-in learner, rather than the model, held \sata back.

\item \textbf{Decisions after data.} Some design choices were made after seeing early data. Most never touched the evaluation tasks: the shortcut's direction was randomised per task after the pilot showed Qwen's ``bigger means 1'' habit backing it; the feature names were picked by their measured effect on Qwen's zero-shot answers on random profiles; and \sata's matched target and drawing rule were chosen on its own validation tasks, judged by the stand-in learner. However, the ten planned tests for \RQ1--\RQ3 had been written down in advance with calibrated balanced accuracy as their measure, and were later switched to \AUROC because, after calibration, a large share of the prompts (about 36\%) gave the same answer to every query.
\end{itemize}

\section{Future work}
The most immediate direct extension is scale, with the same testbed and contrast families on larger models, where the influence of the priors should be weaker. This would show whether selection starts to matter as the model starts to read the evidence. Beyond this, datamodels \citep{ilyas2022datamodels,chang2023data}, where we would train the selector on the model's own scores of demonstration subsets, possibly read at the critical layer, would remove the proxy problem described earlier. One of the parked ideas is to turn the prior rather than fight it, by negating the features whose larger values go with a specific label, in demonstrations and query alike.

